\section{Related work}
\label{sec:related}

HyperPod sits on top of existing work rather than replacing it, and the ideas it
uses are not new. This section records what belongs to others.

\paragraph{\texttt{llama.cpp} and its RPC backend.}
All inference here is performed by \texttt{llama.cpp}: GGUF loading,
quantised kernels, CUDA and Metal execution, and the RPC backend that exposes a
remote accelerator. The ability to combine the memory of several machines is
therefore upstream functionality, not ours. What upstream does not provide, and
explicitly describes as out of scope for a proof-of-concept backend, is
surviving the loss of a peer. That gap is what HyperPod addresses.

\paragraph{Consumer-hardware distributed inference.}
Running a single model across several personal devices has been explored by
other projects, including work that schedules layers according to each device's
capability and work that distributes models across volunteer machines over the
public internet. Capability-aware placement in particular is well established;
our scheduler is a deliberately simple instance of it, using one measured
quantity. We make no novelty claim for the idea of weighting devices by
available memory.

\paragraph{Fault tolerance in serving systems.}
Detecting a failed node, excluding it, and restarting work on the survivors is
standard practice in distributed systems and in datacentre model serving. The
mechanisms used here --- heartbeats, health timeouts, endpoint probing, process
supervision, request retry --- are all conventional. Their combination on three
mismatched laptops over domestic Wi-Fi, with the resulting costs measured, is
the part we contribute.

\paragraph{Meta open-weight models.}
The models are Meta's: Llama~3.2 1B and 3B, and Llama~3.1 8B, used as published
GGUF quantisations. HyperPod does not fine-tune or modify them. Access to
open weights is what makes the work possible at all, because the system moves
model weights between physical devices and reads back where they landed. An
API-only model exposes no such surface.

\paragraph{Where this work differs.}
Three things distinguish what is reported here from the surrounding work. The
placement figures are read out of the inference engine's own load output rather
than assumed from the requested split, which is how we found our tensor-order
bug. Unavailable hardware values are preserved as unavailable all the way to the
interface, so the Mac's missing video-memory figure is visible instead of being
filled in. And the failure experiments were performed by stopping real processes
on real machines during live requests, with the resulting logs committed, rather
than by injecting faults in software.
